\section{早期实现：排序数组后端}
\label{sec:vma-sorted-array}

红黑树是 VMA 后端的终极方案，但它代码量大、调试困难，不适合内核开发的早期阶段。
我们先实现一个\textbf{排序数组}（sorted array）后端——代码不到 150 行，
却足以让 VMA 子系统完整运行，为后续所有章节提供基础。

\begin{designbox}
排序数组的核心权衡：

\begin{itemize}
  \item \textbf{优势}：实现极其简单，不依赖 \texttt{kmalloc}（静态数组），
        \texttt{find} 用二分查找 $O(\log n)$，\texttt{dump} 天然按地址排序
  \item \textbf{劣势}：\texttt{add}/\texttt{remove} 需要搬移数组元素 $O(n)$
  \item \textbf{适用场景}：VMA 数量 $< 64$（内核启动到单进程阶段绰绰有余）
  \item \textbf{升级路径}：Stage 13 引入 per-process VMA 后切换到红黑树后端，
        接口完全兼容，调用者零改动
\end{itemize}
\end{designbox}

\subsection{数据结构}

排序数组直接复用公开的 \texttt{vm\_area\_t}，不需要额外的内部节点类型：

\codefile{src/kernel/mm/vma\_sorted\_array.c}
\begin{lstlisting}[style=cstyle]
#define VMA_MAX_AREAS 256

static vm_area_t g_areas[VMA_MAX_AREAS];
static unsigned  g_count = 0;
\end{lstlisting}

\begin{whybox}
为什么是 256？

每个 \texttt{vm\_area\_t} = 16 字节，$256 \times 16 = 4096$ 字节 = 恰好 1 页。
256 个不重叠区间足以管理 4\,GiB 地址空间（平均每个区间 16\,MiB）。
引导阶段只用 2 个（\texttt{direct-map} + \texttt{kheap}），
VMM 动态分配约 20 个——远不会用完。
\end{whybox}

数组中 \texttt{g\_areas[0..g\_count)} 始终按 \texttt{start} 字段升序排列，
\texttt{g\_count} 跟踪当前已注册的 VMA 数量。

\begin{figure}[H]
  \centering
  % figures/ch11/fig03-vma-sorted-array.tex — auto-extracted from chapter source
\begin{tikzpicture}[
  cell/.style={draw, minimum height=0.8cm, minimum width=1.8cm,
               font=\ttfamily\small, anchor=west},
  lbl/.style={font=\scriptsize\color{gray}, anchor=south},
]
  \node[cell, fill=green!12] (a0) at (0, 0) {area[0]};
  \node[cell, fill=green!12] (a1) at (1.8, 0) {area[1]};
  \node[cell, fill=green!12] (a2) at (3.6, 0) {area[2]};
  \node[cell, fill=gray!8]   (a3) at (5.4, 0) {(空)};
  \node[cell, fill=gray!8]   (a4) at (7.2, 0) {$\cdots$};

  \node[lbl] at (0.9, 0.9) {start=\hex{C000}};
  \node[lbl] at (2.7, 0.9) {start=\hex{D000}};
  \node[lbl] at (4.5, 0.9) {start=\hex{E000}};

  \draw[{Stealth[length=4pt]}-, thick, blue!60!black]
    (5.4, -0.7) -- (5.4, -0.4);
  \node[font=\scriptsize\color{blue!60!black}] at (5.4, -0.9) {\texttt{g\_count = 3}};
\end{tikzpicture}

  \caption{排序数组示意——前 3 个槽位有效，按 \texttt{start} 升序}
  \label{fig:vma-sorted-array}
\end{figure}

与红黑树后端的静态节点池类似，排序数组也完全不依赖 \texttt{kmalloc}——
它是一个编译期 \texttt{static} 数组。这避免了第~\ref{sec:vma-pool}~节讨论的循环依赖问题。

\subsection{二分查找辅助函数}

\texttt{sa\_add}、\texttt{sa\_remove}、\texttt{sa\_find} 都需要在有序数组上做二分查找，
区别仅在于比较条件（\texttt{<} vs \texttt{<=}）。我们抽出两个辅助函数避免重复代码：

\codefile{src/kernel/mm/vma\_sorted\_array.c}
\begin{lstlisting}[style=cstyle]
/* lower_bound: 第一个 start >= target 的下标 */
static unsigned sa_lower_bound(uint32_t target)
{
    unsigned left = 0, right = g_count;
    while (left < right) {
        unsigned mid = left + (right - left) / 2;
        if (g_areas[mid].start < target)
            left = mid + 1;
        else
            right = mid;
    }
    return left;
}

/* upper_bound: 第一个 start > target 的下标 */
static unsigned sa_upper_bound(uint32_t target)
{
    unsigned left = 0, right = g_count;
    while (left < right) {
        unsigned mid = left + (right - left) / 2;
        if (g_areas[mid].start <= target)
            left = mid + 1;
        else
            right = mid;
    }
    return left;
}
\end{lstlisting}

两者的唯一区别是比较时是否包含等号：
\texttt{lower\_bound} 用 \texttt{<}（找 $\geq$），\texttt{upper\_bound} 用 \texttt{<=}（找 $>$）。
\texttt{sa\_add}/\texttt{sa\_remove} 用 \texttt{lower\_bound} 找精确位置，
\texttt{sa\_find} 用 \texttt{upper\_bound} 找"最后一个 $\leq$ addr"的候选。

\subsection{插入：保持有序 + 重叠检测}

\texttt{sa\_add} 的核心逻辑分四步：

\begin{enumerate}
  \item \textbf{参数校验}——容量上限、\texttt{start < end}、页对齐（低 12 位为零）
  \item \textbf{重叠检查}——遍历数组，确认新 \texttt{[start, end)} 不与任何已有 VMA 重叠。
        两个区间重叠的充要条件是 \texttt{a.start < b.end \&\& b.start < a.end}
  \item \textbf{找插入位置}——因数组有序，找到第一个 \texttt{g\_areas[i].start > start} 的位置
  \item \textbf{后移 + 写入}——从尾部往回搬，腾出插入位，写入新 VMA
\end{enumerate}

\begin{figure}[H]
  \centering
  % figures/ch11/fig04-vma-sa-insert.tex — auto-extracted from chapter source
\begin{tikzpicture}[
  cell/.style={draw, minimum height=0.7cm, minimum width=1.2cm,
               font=\ttfamily\small, anchor=west},
  highlight/.style={fill=yellow!30},
  arr/.style={-{Stealth[length=4pt]}, thick, red!60!black},
]
  % 插入前
  \node[font=\small\bfseries] at (-1.5, 0) {插入前:};
  \node[cell, fill=green!12] at (0, 0) {A};
  \node[cell, fill=green!12] at (1.2, 0) {B};
  \node[cell, fill=green!12] at (2.4, 0) {D};
  \node[cell, fill=green!12] at (3.6, 0) {E};
  \node[cell, fill=gray!8]   at (4.8, 0) {\_};

  % 后移
  \node[font=\small\bfseries] at (-1.5, -1.2) {后移:};
  \node[cell, fill=green!12] at (0, -1.2) {A};
  \node[cell, fill=green!12] at (1.2, -1.2) {B};
  \node[cell, highlight]     at (2.4, -1.2) {\_};
  \node[cell, fill=green!12] at (3.6, -1.2) {D};
  \node[cell, fill=green!12] at (4.8, -1.2) {E};
  \draw[arr] (3.15, 0.35) -- (4.35, -0.85);
  \draw[arr] (4.35, 0.35) -- (5.55, -0.85);

  % 写入
  \node[font=\small\bfseries] at (-1.5, -2.4) {写入:};
  \node[cell, fill=green!12] at (0, -2.4) {A};
  \node[cell, fill=green!12] at (1.2, -2.4) {B};
  \node[cell, fill=red!15]   at (2.4, -2.4) {C};
  \node[cell, fill=green!12] at (3.6, -2.4) {D};
  \node[cell, fill=green!12] at (4.8, -2.4) {E};
  \node[font=\scriptsize\color{red!60!black}] at (3.0, -3.1) {新插入};
\end{tikzpicture}

  \caption{\texttt{sa\_add} 插入过程：后移元素 → 写入新 VMA}
  \label{fig:vma-sa-insert}
\end{figure}

\subsection{查找：二分搜索}

\texttt{sa\_find(addr)} 调用 \texttt{sa\_upper\_bound} 找到第一个 \texttt{start > addr}
的位置，减 1 即为候选者：

\codefile{src/kernel/mm/vma\_sorted\_array.c}
\begin{lstlisting}[style=cstyle]
static const vm_area_t* sa_find(uint32_t addr)
{
    if (g_count == 0)
        return NULL;

    unsigned idx = sa_upper_bound(addr);
    if (idx == 0)
        return NULL;

    const vm_area_t* candidate = &g_areas[idx - 1];
    return (addr < candidate->end) ? candidate : NULL;
}
\end{lstlisting}

算法步骤：

\begin{enumerate}
  \item 初始化 \texttt{left = 0}, \texttt{right = g\_count}
  \item 当 \texttt{left < right} 时，计算 \texttt{mid = left + (right - left) / 2}
  \item 如果 \texttt{g\_areas[mid].start <= addr}，则 \texttt{left = mid + 1}
        （\texttt{mid} 可能是答案，但看右边有没有更后面的候选者）
  \item 否则 \texttt{right = mid}（\texttt{mid} 的 \texttt{start} 太大，答案在左边）
  \item 循环结束后，\texttt{left - 1} 是最后一个 \texttt{start <= addr} 的下标
  \item 检查 \texttt{addr < candidate->end}，是则命中，否则落在空洞中
\end{enumerate}

\begin{whybox}
为什么用 upper\_bound 而不是 lower\_bound？

我们要找的是"最后一个 \texttt{start <= addr} 的 VMA"。
\texttt{upper\_bound} 找的是第一个 \texttt{start > addr} 的位置，
减 1 就是我们要的答案。
这是二分查找最常见的变体之一，值得牢记。

如果 \texttt{left == 0}（所有 VMA 的 \texttt{start} 都大于 \texttt{addr}），
直接返回 NULL——不存在包含该地址的 VMA。
\end{whybox}

\begin{figure}[H]
  \centering
  % figures/ch11/fig05-vma-sa-find.tex — auto-extracted from chapter source
\begin{tikzpicture}[
  cell/.style={draw, minimum height=0.8cm, minimum width=2.8cm,
               font=\ttfamily\scriptsize, anchor=west},
  hit/.style={fill=green!20},
  miss/.style={fill=red!10},
]
  \node[cell, fill=blue!8] (a0) at (0, 0) {[0x1000, 0x5000)};
  \node[cell, fill=blue!8] (a1) at (2.8, 0) {[0x8000, 0xA000)};
  \node[cell, fill=blue!8] (a2) at (5.6, 0) {[0xC000, 0xF000)};

  % find(0x9000) → hit
  \draw[-{Stealth[length=4pt]}, thick, green!50!black]
    (4.2, 1.0) -- (4.2, 0.4);
  \node[font=\scriptsize\color{green!50!black}] at (4.2, 1.2)
    {find(0x9000) → area[1] ✓};

  % find(0xB000) → miss (gap)
  \draw[-{Stealth[length=4pt]}, thick, red!50!black]
    (5.0, -1.0) -- (5.0, -0.4);
  \node[font=\scriptsize\color{red!50!black}] at (5.0, -1.2)
    {find(0xB000) → NULL (空洞)};
\end{tikzpicture}

  \caption{二分查找示例：0x9000 命中 area[1]；0xB000 落在两个 VMA 之间的空洞}
  \label{fig:vma-sa-find}
\end{figure}

\subsection{删除：前移覆盖}

\texttt{sa\_remove(start)} 用 \texttt{sa\_lower\_bound} 二分定位后，将后续元素前移一格覆盖：

\begin{figure}[H]
  \centering
  % figures/ch11/fig06-vma-sa-remove.tex — auto-extracted from chapter source
\begin{tikzpicture}[
  cell/.style={draw, minimum height=0.7cm, minimum width=1.2cm,
               font=\ttfamily\small, anchor=west},
  arr/.style={-{Stealth[length=4pt]}, thick, blue!60!black},
]
  \node[font=\small\bfseries] at (-1.5, 0) {删除前:};
  \node[cell, fill=green!12] at (0, 0) {A};
  \node[cell, fill=green!12] at (1.2, 0) {B};
  \node[cell, fill=red!15]   at (2.4, 0) {C};
  \node[cell, fill=green!12] at (3.6, 0) {D};
  \node[cell, fill=green!12] at (4.8, 0) {E};

  \node[font=\small\bfseries] at (-1.5, -1.2) {前移:};
  \node[cell, fill=green!12] at (0, -1.2) {A};
  \node[cell, fill=green!12] at (1.2, -1.2) {B};
  \node[cell, fill=green!12] at (2.4, -1.2) {D};
  \node[cell, fill=green!12] at (3.6, -1.2) {E};
  \node[cell, fill=gray!8]   at (4.8, -1.2) {\_};
  \draw[arr] (4.35, 0.35) -- (3.15, -0.85);
  \draw[arr] (3.15, 0.35) -- (1.95, -0.85);

  \node[font=\scriptsize\color{red!60!black}] at (3.0, 0.55) {删除 C};
\end{tikzpicture}

  \caption{\texttt{sa\_remove} 删除过程：前移元素覆盖被删除位置}
  \label{fig:vma-sa-remove}
\end{figure}

\subsection{复杂度总结与适用阶段}

\begin{table}[H]
\centering
\begin{tabular}{lcc}
\toprule
\textbf{操作} & \textbf{排序数组} & \textbf{适用场景} \\
\midrule
\texttt{find}   & $O(\log n)$ 二分    & 每次 \#PF 都要调用，必须快 \\
\texttt{add}    & $O(n)$ 搬移        & 仅启动/分配时调用，频率低 \\
\texttt{remove} & $O(n)$ 搬移        & 仅释放时调用，频率低 \\
\texttt{dump}   & $O(n)$ 顺序遍历    & 仅调试命令，不影响性能 \\
\bottomrule
\end{tabular}
\caption{排序数组后端复杂度——\texttt{find} 快，\texttt{add}/\texttt{remove} 可接受}
\label{tab:vma-sa-complexity}
\end{table}

在内核启动到单进程阶段（Stage 9 -- Stage 12），VMA 数量 $n < 20$。
此时 $O(n)$ 的插入/删除在绝对耗时上几乎无开销（搬移 20 个 16 字节的结构体 $\approx$ 几十个 \texttt{mov} 指令）。
当多进程引入 per-process VMA 后再切换红黑树，是更合理的渐进式开发策略。

% ============================================================================

